%! Convolutional Neural Nets — Unit 8, Lesson 8.2 — Slides (Beamer)
\documentclass[aspectratio=169, 11pt]{beamer}
\usepackage{linalg-beamer}   % provides \burgundyheader{} and \sectionlabel[color]{}
\newcommand{\vv}{\mathbf{v}}
\newcommand{\bb}{\mathbf{b}}
\newcommand{\ww}{\mathbf{w}}
\newcommand{\relu}{\mathrm{ReLU}}

\begin{document}

% TITLE SLIDE (CourseName is not defined in beamer, so the name is literal)
{
\setbeamercolor{background canvas}{bg=burgundy}
\begin{frame}[plain]
  \vspace{2.8cm}\hspace{0.55cm}%
  \begin{minipage}{0.9\textwidth}
    {\color{goldacc}\bfseries\small LESSON 8.2}\\[0.5cm]
    {\color{white}\bfseries\LARGE Convolutional Neural Nets}\\[1.0cm]
    {\color{blushmid}\small Linear Algebra~$\cdot$~Unit 8: Learning from Data}
  \end{minipage}
\end{frame}
}

% HOOK
\begin{frame}[t, plain]
  \burgundyheader{Too many weights --- and the pattern repeats}
  \sectionlabel{Hook}
  \begin{columns}[T]
    \begin{column}{0.54\textwidth}
      A \textbf{full} layer $\relu(A\vv+\bb)$ gives \emph{every} connection its own weight.\\[6pt]
      A photo has \emph{millions} of pixels --- that is $\sim\!10^{12}$ weights. And a bright \textbf{edge}
      looks the same in the corner or the center.\\[6pt]
      \textbf{Idea:} learn \emph{one} small pattern-detector and \textbf{reuse} it everywhere.
    \end{column}
    \begin{column}{0.46\textwidth}
      \centering
      \begin{tikzpicture}[scale=0.5]
        \draw[->, charcoal, thin] (-0.3,0)--(5.4,0) node[right]{\scriptsize position};
        \draw[->, charcoal, thin] (0,-0.4)--(0,3.2) node[above]{\scriptsize value};
        \draw[ultra thick, royalblue] (1,0.8)--(2,0.8)--(3,0.8)--(4,2.4)--(5,2.4);
        \foreach \x/\y in {1/0.8,2/0.8,3/0.8,4/2.4,5/2.4}
          \fill[royalblue] (\x,\y) circle (1.6pt);
        \draw[goldacc, very thick, dashed] (2.75,0.45)--(4.25,0.45)--(4.25,2.8)--(2.75,2.8)--cycle;
        \node[goldacc] at (3.5,3.15) {\scriptsize filter slides};
      \end{tikzpicture}
    \end{column}
  \end{columns}
  \vspace{2pt}
  \begin{block}{The question of Lesson 8.2}
    How can one small filter, reused everywhere, detect a pattern \textbf{wherever} it appears --- with almost
    no weights?
  \end{block}
\end{frame}

% SLIDE THE FILTER
\begin{frame}[t, plain]
  \burgundyheader{Convolution: slide a small filter}
  \sectionlabel[goldacc]{Skill}
  \begin{columns}[T]
    \begin{column}{0.5\textwidth}
      Filter $\ww=[-1,1]$, signal $\vv=[2,2,2,6,6]$. At each spot: $-v_i+v_{i+1}$ (``next minus current'').
      \[ [\,0,\ 0,\ 4,\ 0\,] \]
      Zero on the flat parts, a \textbf{spike} at the edge $2\to6$. This output is the \textbf{feature map} ---
      big where the pattern is.
    \end{column}
    \begin{column}{0.5\textwidth}
      \centering
      \begin{tikzpicture}[scale=0.5]
        \draw[->, charcoal, thin] (-0.3,0)--(4.8,0) node[right]{\scriptsize position};
        \draw[->, charcoal, thin] (0,-0.4)--(0,3.2) node[above]{\scriptsize output};
        \foreach \x/\y in {1/0,2/0,3/1.7,4/0}{
          \draw[ultra thick, burgundy] (\x,0)--(\x,\y);
          \fill[burgundy] (\x,\y) circle (1.6pt);
        }
        \node[burgundy] at (3.0,2.15) {\scriptsize $4$};
        \node[charcoal] at (3.0,-0.8) {\scriptsize fires at the edge};
      \end{tikzpicture}
    \end{column}
  \end{columns}
\end{frame}

% IT IS A MATRIX
\begin{frame}[t, plain]
  \burgundyheader{It is a matrix --- with shared weights}
  \sectionlabel{Skill}
  \[ A=\begin{bmatrix} -1 & 1 & 0 & 0 & 0\\ 0 & -1 & 1 & 0 & 0\\ 0 & 0 & -1 & 1 & 0\\ 0 & 0 & 0 & -1 & 1 \end{bmatrix},
     \qquad A\,[2,2,2,6,6]^{\!\top}=[0,0,4,0]^{\!\top}. \]
  \begin{itemize}
    \setlength{\itemsep}{5pt}
    \item The same $-1,1$ \textbf{repeat down the diagonals} --- a \textbf{convolution matrix}.
    \item \textbf{Weight count:} a full $4\times5$ block $=20$ weights; this convolution $=\mathbf{2}$.
    \item A $1000\times1000$ image: $\sim\!10^{12}$ full weights vs.\ a $3\times3$ filter's $\mathbf{9}$ ---
          reused everywhere. That is \textbf{weight sharing}.
  \end{itemize}
\end{frame}

% CONV LAYER + WHY IT MATTERS
\begin{frame}[t, plain]
  \burgundyheader{A convolutional layer --- and the road ahead}
  \sectionlabel[goldacc]{Payoff}
  \begin{itemize}
    \setlength{\itemsep}{6pt}
    \item Same layer as before: $\relu(A\vv+\bb)$, only now $A$ is a \textbf{convolution matrix} and $\bb$ is a
          \emph{threshold}. On $[0,0,4,0]$ with bias $-1$: $\relu([-1,-1,3,-1])=[0,0,3,0]$ --- keep the strong response.
    \item \textbf{Detect anywhere:} one reused filter fires on its pattern \emph{wherever} it sits --- a cat in
          the corner or the center (translation invariance).
    \item \textbf{Why it matters:} weight sharing turns a trillion weights into a handful --- small enough to
          actually learn. \emph{Still just matrices $+$ ReLU.}
  \end{itemize}
  \vspace{2pt}
  \begin{block}{The road ahead}
    \textbf{8.3} minimizing loss by gradient descent \;$\bullet$\; \textbf{8.4} mean, variance, covariance.
  \end{block}
\end{frame}

\end{document}
